% Source: https://artofproblemsolving.com/community/c3h2159242_wheels__circles
% Stable identity: aops-topic-2159242-q01
In one revolution a wheel travels its circumference: $20\pi$ inches for the larger wheel and $17\pi$ inches for the smaller one, a difference of $3\pi$ inches per revolution.

\begin{asy}
usepackage("amsmath");
size(6cm);
defaultpen(fontsize(10pt));
draw(circle((0,10),10),linewidth(1)); draw(circle((30,8.5),8.5),linewidth(1));
draw((-12,0)--(45,0),linewidth(0.8));
draw((-10,10)--(10,10),blue); draw((21.5,8.5)--(38.5,8.5),blue);
label("$20$ in",(0,10),N,blue); label("$17$ in",(30,8.5),N,blue);
label("circumference $20\pi$",(0,0),S); label("circumference $17\pi$",(30,0),S);
label("\mbox{}",(min(currentpicture,user=true)+max(currentpicture,user=true))/2,black); // keeps the text after the figure black when LaTeX typesets the labels inline
\end{asy}

In $30$ minutes $=1800$ seconds each wheel makes $10\cdot1800=18000$ revolutions. So the larger wheel travels
\[18000\cdot3\pi=\boxed{54000\pi\text{ inches}}=4500\pi\text{ feet}\approx14137\text{ feet}\]
farther (about $2.68$ miles).
